% =====================================================================
\section{Experiments}\label{sec:exp}
% =====================================================================

We evaluate Walk Attention on a controlled family of graphs in which
over-squashing is severe and \emph{tunable}, so that the contribution of each
architectural ingredient can be isolated. All code, configurations and a full
experiment-tracking log are released with the paper.

% ---------------------------------------------------------------------
\subsection{A tunable over-squashing benchmark}
% ---------------------------------------------------------------------

\paragraph{Bottleneck-chain graphs.}
We build a family parameterised by a \emph{depth} $d$, a number of sources $K$
and a width $M$. The graph has $K$ source vertices, then $d$ consecutive
\emph{bottleneck layers} of $M$ interchangeable vertices each, and finally one
target vertex. Every source connects to every vertex of the first layer, every
vertex of layer $r$ connects to every vertex of layer $r+1$, and every vertex of
the last layer connects to the target. Because the bottleneck vertices are
interchangeable, the number of length-$(d{+}1)$ walks from a source to the target
is $M^{d}$, so the walk multiplicity into the target grows as
$n_{d+1}(\cdot,\text{target})\sim K\,M^{d}$
(Proposition~\ref{prop:count}); a message-passing GNN must squash all of them
into the target's fixed-width vector. The walk-reachability support, by contrast,
remains sparse: at the deepest range it is only the $K$ source--target pairs,
about $2\%$ of all node pairs in our configuration.

\paragraph{Task.}
Each source carries a one-hot \emph{key} and a scalar \emph{payload}; the target
carries a query key. The model must output, at the target node, the key of the
source whose key matches the query---a retrieval task whose answer sits $d{+}1$
hops away, behind the bottleneck. Accuracy is measured at the target node;
chance is $1/K$.

\paragraph{Models.} We compare:
\begin{itemize}
  \item \textbf{GAT} \cite{velickovic2017graph}: one-hop graph attention, $d{+}1$
  layers (learned weights, one-hop support);
  \item \textbf{Walk operator} (\texttt{walkraw}): the fixed-weight multi-hop
  operator $W^{(g)}$ of Definition~\ref{def:walkop} (multi-hop support, fixed
  multiplicity weights);
  \item \textbf{Walk Attention} (ours): Algorithm~\ref{alg:wa} (multi-hop
  support, learned weights).
\end{itemize}
All models use hidden width $16$, four heads where applicable, and $d{+}1$
layers/ranges. We train with AdamW, a ten-epoch linear learning-rate warmup
followed by cosine decay, and gradient clipping; Walk Attention additionally uses
layer normalisation after each layer. We report mean and standard deviation of
target accuracy over five random seeds, at problem depths $d\in\{2,3\}$ with
$K=5$, $M=4$.

% ---------------------------------------------------------------------
\subsection{Main result}
% ---------------------------------------------------------------------

Table~\ref{tab:main} reports the results. One-hop GAT collapses towards chance
($0.26$--$0.44$): it cannot carry the signal across the bottleneck, as expected
of a local operator under over-squashing. The fixed-weight walk operator reaches
the target---its multi-hop support bypasses the bottleneck---but, weighting all
reachable sources by multiplicity alone, it cannot select the matching source and
plateaus at $0.50$--$0.62$. \textbf{Walk Attention solves the task exactly},
attaining $1.000\pm0.000$ accuracy at both depths across all five seeds: the same
multi-hop support, now with learned query--key weights, lets the target attend
selectively to the relevant source.

\begin{table}[t]
\centering
\caption{Target accuracy (mean $\pm$ std over 5 seeds) on bottleneck-chain
retrieval, $K{=}5$, $M{=}4$. Chance is $1/K=0.20$. Depth $d$ is the number of
bottleneck layers; the answer lies $d{+}1$ hops away.}
\label{tab:main}
\begin{tabular}{lcc}
\toprule
\textbf{Model} & \textbf{depth $d=2$} & \textbf{depth $d=3$} \\
\midrule
GAT (one-hop attention)            & $0.44\pm0.12$ & $0.26\pm0.10$ \\
Walk operator (fixed weights)      & $0.62\pm0.12$ & $0.50\pm0.12$ \\
\textbf{Walk Attention (ours)}     & $\mathbf{1.000\pm0.000}$ & $\mathbf{1.000\pm0.000}$ \\
\bottomrule
\end{tabular}
\end{table}

\paragraph{On stability.}
Without layer normalisation and warmup, Walk Attention has a high ceiling but is
high-variance: across seeds its accuracy ranged from $0.44$ to $1.00$. The three
standard stabilisers (layer normalisation, learning-rate warmup, gradient
clipping) remove this variance entirely, yielding the $\pm0.000$ in
Table~\ref{tab:main}. The instability was therefore an optimisation artefact, not
a property of the operator.

% ---------------------------------------------------------------------
\subsection{A negative result: collapsing walks hurts}
% ---------------------------------------------------------------------

The path-algebra quotient $\kQ/I$ of Section~\ref{sec:algebra} suggests an
alternative remedy: replace the walk multiplicities $n_g(i,j)$ by the smaller
quotient dimensions $\dim(e_j(\kQ/I)_g e_i)$, de-duplicating equivalent walks
before aggregation. We implemented this (the fixed-weight operator built from the
quotient dimensions) and compared it to the unquotiented walk operator on the
same task, in a separate controlled ablation that changes \emph{only} the
operator (quotient vs.\ raw) and holds everything else fixed. As
Table~\ref{tab:quotient} shows, \textbf{the quotient operator underperforms the
raw one}: collapsing parallel walks discards the source multiplicity that the
retrieval task needs. In this regime redundant paths are \emph{signal}, not
noise, and removing them removes information.

\begin{table}[t]
\centering
\caption{De-duplicating walks via the quotient $\kQ/I$ \emph{hurts}: target
accuracy (mean over 3 seeds), bottleneck retrieval. This is a separate ablation
(3 seeds, without the stabilisers of Table~\ref{tab:main}) that isolates the
operator change, so the walk-operator column is not directly comparable to
Table~\ref{tab:main}; what matters here is the within-row gap: the quotient
operator is consistently below the unquotiented one.}
\label{tab:quotient}
\begin{tabular}{lcc}
\toprule
\textbf{Operator} & \textbf{depth $d=2$} & \textbf{depth $d=3$} \\
\midrule
Quotient $\kQ/I$ (collapsed walks) & $0.47$ & $0.52$ \\
Walk operator (raw walks)          & $\mathbf{0.64}$ & $\mathbf{0.53}$ \\
\bottomrule
\end{tabular}
\end{table}

This negative result is informative: it tells us that the value of the path
algebra here is \emph{not} compression. The right use of the algebra is to supply
the multi-hop attention \emph{support} (which it does, sparsely and exactly), and
to let the network learn weights over that support---which is precisely Walk
Attention. We regard finding a regime in which the quotient's compression is
itself beneficial (e.g.\ when parallel paths are genuinely redundant noise, or
when distinct path classes should be routed to distinct heads) as an open
question rather than a settled negative.

% ---------------------------------------------------------------------
\subsection{A real long-range benchmark: LRGB Peptides-func}
% ---------------------------------------------------------------------

The synthetic benchmark is deliberately a worst case: an extreme bottleneck where
multi-hop reach is decisive. To see how Walk Attention behaves on real data we
also evaluate on \textbf{Peptides-func} from the Long Range Graph Benchmark
\cite{dwivedi2022lrgb} ($10{,}873/2{,}331/2{,}331$ train/val/test molecular
graphs, median $138$ nodes; multi-label graph classification, metric Average
Precision). We use the same atom embedding, four layers, hidden width $80$, mean
pooling and AdamW for every model; Walk Attention uses ranges $g=1,\dots,4$ with
the walk-reachability support of each molecule. On these graphs the support stays
sparse ($1$--$8\%$ of $n^2$ at range $3$), so the method remains cheap.

Table~\ref{tab:lrgb} reports test AP. \textbf{Walk Attention is competitive with
the strong baselines but does not surpass them here}: it clearly beats GCN and is
statistically indistinguishable from GAT. This is informative and we report it
plainly. Peptides graphs are small and chain-like, without the acute bottlenecks
of the synthetic family, so a few layers of one-hop attention already reach far
enough; the multi-hop support that is decisive under severe over-squashing buys
little when over-squashing is mild. Walk Attention's advantage is therefore
\emph{regime-dependent}: it is large exactly when the squashing is severe
(Table~\ref{tab:main}) and shrinks to a tie when it is not. Crucially, on real
data the operator still trains stably and remains cheap; it does not degrade.

\begin{table}[t]
\centering
\caption{LRGB Peptides-func, test Average Precision (higher is better). Walk
Attention is competitive with GAT and above GCN; on this mild-over-squashing
benchmark the multi-hop support yields no decisive gain (contrast
Table~\ref{tab:main}).}
\label{tab:lrgb}
\begin{tabular}{lc}
\toprule
\textbf{Model} & \textbf{test AP} \\
\midrule
GCN                          & $0.484$ \\
GAT                          & $\mathbf{0.527}$ \\
Walk Attention (ours)        & $0.522$ \\
\bottomrule
\end{tabular}
\end{table}

% ---------------------------------------------------------------------
\subsection{Cost}
% ---------------------------------------------------------------------

Walk Attention operates on the walk-reachability support, whose size is the
number of nonzero entries of $\Adj^{\,g}$. On the bottleneck family this is a
small fraction of $n^2$ (about $2\%$ at the deepest range), so the layer is far
cheaper than a dense graph Transformer while retaining its long range. The
supports are precomputed once per graph topology and cached, so the path-algebra
computation does not enter the training loop.
